\clearpage
\pdfbookmark[0]{Preface.}{praefatio}
\begin{center}\Large PREFACE.\end{center}
\medskip
\SourceLeaf{010}Once \textsc{Kepler}'s genius had discovered the laws of planetary motion, there was no lack of means for deriving the elements of each planet from observations. \textsc{Tycho Brahe}, who had raised practical astronomy to a height previously unknown, had observed all the planets for a long succession of years with such care and perseverance that \textsc{Kepler}, the most worthy heir to such a treasure, had only to select the observations that seemed suited to whatever purpose he had in view. The work was also considerably lightened by the mean motions of the planets, which had long since been determined with the greatest precision from the oldest observations.

The astronomers who after \textsc{Kepler} sought to measure the orbits of the planets still more accurately with the aid of more recent or more precise observations enjoyed the same or even greater advantages. There were no longer any wholly unknown elements to determine. The known elements merely needed slight correction and had to be confined within narrower limits.

The principle of universal gravitation discovered by the great \textsc{Newton} opened up a wholly new field. It showed that, with only a slight modification, the same laws that \textsc{Kepler} had found to govern the five planets must necessarily
\SourceLeaf{011}be obeyed by \emph{all} celestial bodies whose motions are governed solely by the force of the Sun. Relying on the evidence of observations, \textsc{Kepler} had declared that the orbit of every planet is an ellipse in which areas are swept out uniformly about the Sun, which occupies one of the foci, and that the periods of revolution in different ellipses are in the ratio of the \(3/2\) powers of their major semiaxes. Starting from the principle of universal gravitation, \textsc{Newton} demonstrated a priori that all bodies governed by the attractive force of the Sun must move in conic sections. The planets show us one kind, namely ellipses. The other kinds, parabolas and hyperbolas, must be regarded as equally possible, provided there are bodies moving at the requisite speeds under the force of the Sun. He also showed that the Sun always occupies one focus of the conic section, that the areas swept out about the Sun by a given body in different intervals are proportional to those intervals, and finally that the areas swept out about the Sun by different bodies in equal intervals are in the ratio of the square roots of the semiparameters of their orbits. The last of these laws, which for elliptical motion is identical to \textsc{Kepler}'s last law, extends to parabolic and hyperbolic motion, to which Kepler's law cannot apply because there are no revolutions. A thread had now been found to guide the way into the previously inaccessible labyrinth of cometary motions. This succeeded so well that a single hypothesis, that the orbits were parabolas, sufficed to account for the motions of all comets that had been observed accurately. The system of universal gravitation had thus secured new and most splendid triumphs for analysis. Comets, which until then had always been untamed or, when they seemed vanquished, had soon become unruly and rebellious again, submitted to the reins. Enemies turned into guests, they continued on their way along the paths traced by calculation, faithfully obeying the same eternal laws as the planets.

Even the determination of parabolic cometary orbits from observations presented far greater difficulties than the determination of the elliptical orbits of planets. This was chiefly because comets were visible for a shorter
\SourceLeaf{012}interval, so the mathematician could not choose observations particularly suited to one purpose or another. He was forced to use those that chance had provided, and the special methods used in calculating planetary orbits could scarcely ever be applied. The great \textsc{Newton} himself, the foremost mathematician of his century, made no secret of the difficulty of the problem, yet, as might be expected, he emerged victorious from this contest too. Many mathematicians after \textsc{Newton} applied themselves to the same problem, with varying success, and so little now remains to be desired.

One should bear in mind, however, how conveniently the difficulty of this problem too is reduced by knowledge of one element of the conic section, since the very assumption that the orbit is parabolic makes the major axis infinitely large. If their positions are disregarded, all parabolas differ from one another solely by the greater or lesser distance of the vertex from the focus, whereas conic sections in general admit infinitely greater variety. There was certainly no sufficient reason to suppose that the trajectories of comets were parabolic with absolute precision. Indeed, it must be considered infinitely improbable that nature should ever have conformed to such an assumption. It was known, however, that near perihelion the phenomena of a celestial body moving in an ellipse or hyperbola whose major axis is very large relative to its parameter differ very little from motion in a parabola with the same distance from vertex to focus, and that the difference becomes smaller as the ratio of axis to parameter increases. Experience had also shown that the residual differences between the observed motion and the motion calculated in a parabolic orbit were scarcely ever greater than could safely be attributed to the errors in the observations themselves, which in this case were usually quite appreciable. Astronomers therefore thought it best to stop at the parabola. They were quite right, since there was no means at all of determining with sufficient certainty whether there were any departures from the parabola or how large they might be. An exception must be made for the famous comet of \textsc{Halley}, which describes a very elongated ellipse and revealed its period to us through observations made on several of its
\SourceLeaf{013}returns to perihelion. Once the major axis is known from the period, the calculation of the remaining elements can scarcely be considered more difficult than the determination of a parabolic orbit. It should also be mentioned that for several other comets observed over a somewhat longer time, astronomers tried to determine the departure from a parabola. All the methods proposed or used for this purpose, however, rest on the assumption that the departure from the parabola is small. In these attempts, the parabola already calculated therefore supplied approximate values for the individual elements, apart from the major axis or the period of revolution that depends on it. These values required only slight correction. It must also be admitted that all these attempts scarcely ever yielded a definite result, with the possible exception of the comet of 1770.

As soon as it was recognized that the motion of the new planet discovered in 1781 could not be reconciled with the parabolic hypothesis, astronomers began to fit a circular orbit to it, a task accomplished by a very easy and simple calculation. As luck would have it, the orbit of this planet was only moderately eccentric, so the elements derived ·from that assumption at least supplied a rough approximation on which the determination of the elliptical elements could then be based. Several other favourable circumstances came into play. The slow motion of the planet and the very small inclination of its orbit to the plane of the ecliptic made the calculations much simpler and made it possible to employ special methods unsuitable in other cases. They also dispelled the fear that the planet, immersed in the Sun's rays, might later elude those who searched for it. Under other circumstances, especially if its light had also been fainter, that fear could certainly have caused concern. A more precise determination of the orbit could therefore safely be postponed until it was possible to choose, from more numerous and more widely separated observations, those that seemed best suited to the purpose.

\SourceLeaf{014}Thus, in all cases in which the orbits of celestial bodies had to be deduced from observations, there were appreciable advantages that suggested or at least permitted the use of special methods. The chief advantage was that hypothetical assumptions had already made it possible to obtain approximate values of some elements before the calculation of the elliptical elements was undertaken. It nevertheless seems rather surprising that the general problem

\begin{quote}\itshape
To determine the orbit of a celestial body, without any hypothetical assumption, from observations covering a short time and consequently allowing no selection suited to the application of special methods
\end{quote}
\noindent was almost entirely neglected until the beginning of this century, or at least was never given serious and adequate treatment. Its difficulty and elegance could certainly have commended it to theorists, even if its great usefulness was still unrecognized by practitioners. The prevailing opinion was that such a complete determination from observations confined to a short interval was \emph{impossible}. This opinion was surely ill-founded, since it has now been established with the greatest certainty that the orbit of a celestial body can already be determined to a sufficient approximation, without any hypothetical assumption, from good observations covering only a few days.

In September 1801, while engaged in quite different work, I had come upon a few ideas that seemed suited to solving the great problem I have just described. In such a situation, to keep from being distracted too much from an investigation we find absorbing, we often let associations of ideas slip away unnoticed, though on closer examination they might have borne the richest fruit. Perhaps the same fate awaited those little ideas of mine, had they not come at a time more favourable to their preservation and development than any one could have chosen. At about that same time, word was on everyone's lips that a new planet had been discovered on 1 January of that year at the Palermo observatory. Soon the observations made by the distinguished astronomer \textsc{Piazzi} from
\SourceLeaf{015}that date to 11 February became public. Surely nowhere in the annals of astronomy do we find so momentous an occasion to demonstrate the importance of this problem as clearly as we could then, in so critical a situation and under such urgent necessity. One could scarcely imagine a more momentous one. All hope of finding that planetary atom again in the sky among countless tiny stars, after nearly a year had elapsed, depended entirely on a sufficiently close approximation to its orbit based on those few observations alone. Could I ever have had a better opportunity to test what my little ideas were worth in practice than by using them then to determine the orbit of Ceres? In those 41 days the planet had traced a geocentric arc of only three degrees, and after a year had elapsed it had to be sought in a region of the sky very far from there. This first application of the method took place in October 1801. The very first clear night on which the planet was sought using the numbers thus obtained\footnote{7 December 1801, by the distinguished de \textsc{Zach}.} returned the fugitive to observation. Three other new planets discovered since then have provided fresh occasions to test and confirm the effectiveness and generality of the method.

Immediately after the recovery of Ceres, several astronomers wished me to publish the methods I had used for those calculations. Several things kept me from complying with these friendly requests at the time. I had other occupations, I wanted to treat the subject more fully at some point, and above all I expected that continued work on the investigation would bring the various parts of the solution to a higher degree of generality, simplicity, and elegance. Since this hope has been fulfilled, I think I have no reason to regret the delay. The methods first employed underwent so many changes, and such great ones, that scarcely a trace of resemblance remains between the way the orbit of Ceres was originally calculated and the procedure
\SourceLeaf{016}set out in this work. Although it was beyond my purpose to write a full account of all these investigations as they gradually took more and more perfect form, on several occasions, especially whenever a problem of some importance was at issue, I thought that some account of the earlier methods should also be retained. In addition to the solutions of the principal problems, I have treated in this work many matters that came to my attention during a fairly long study of the motions of celestial bodies in conic sections and seemed especially worth attending to for their analytical elegance or, above all, for their practical use. I have always given greater care to matters or methods of my own, touching only lightly on familiar ones and only as far as the connection of the subject seemed to require.

The whole work is accordingly divided into two parts. The first book develops the relations between the quantities on which the motions of celestial bodies about the Sun depend according to \textsc{Kepler}'s laws. The first two sections develop the relations for a single position considered on its own, and the third and fourth sections develop those for the comparison of several positions. The first two sections present both the methods in common use and, above all, others that I believe are far preferable in practice, by which one passes from known elements to the phenomena. The latter two deal with many highly important problems that pave the way for the inverse operations. The phenomena themselves are formed by a subtle and intricate interweaving of the elements. The way they are woven together must be understood more thoroughly before one can hope to disentangle the threads and resolve the work into its elements. The first book thus furnishes the instruments and resources with which this arduous task itself is then carried out in the second book. The greater part of the labour then consists in assembling those resources properly, arranging them in a suitable order, and directing them toward the proposed goal.

Most of the more important problems have been illustrated with suitable
\SourceLeaf{017}examples, drawn whenever possible from actual observations. In this way readers will gain greater confidence in the effectiveness of the methods and see more clearly how to use them. I also hope that even those with less experience will be encouraged to study matters that undoubtedly constitute the most fruitful and beautiful part of theoretical astronomy.

I wrote this at Göttingen on 28 March 1809.
